```latex
\documentclass[11pt]{article}
\usepackage{amsmath,amssymb,amsthm}
\usepackage{hyperref}

\newtheorem{theorem}{Theorem}
\newtheorem{lemma}{Lemma}
\newtheorem{proposition}{Proposition}
\theoremstyle{remark}
\newtheorem{remark}{Remark}

\DeclareMathOperator{\Teich}{Teich}
\DeclareMathOperator{\Mod}{Mod}
\DeclareMathOperator{\Def}{Def}
\DeclareMathOperator{\Stab}{Stab}

\title{An explicit Dehn twist with infinite orbit on the components of a deformation space}
\author{}
\date{}

\begin{document}
\maketitle

Let \(f:\mathbb{S}^2\to \mathbb{S}^2\) be a quadratic rational map with critical points \(c_1,c_2\).
Assume that \(c_1\) belongs to a cycle of length \(4\), and that the critical value
\(v=f(c_2)\) does not belong to that cycle. Write
\[
a_0=c_1,\qquad a_1=f(c_1),\qquad a_2=f^{\circ 2}(c_1),\qquad a_3=f^{\circ 3}(c_1),
\]
so that
\[
f(a_0)=a_1,\quad f(a_1)=a_2,\quad f(a_2)=a_3,\quad f(a_3)=a_0.
\]
Set
\[
A=\{a_0,a_1,a_2,a_3\},\qquad B=A\cup\{v\}.
\]
Since \(v\notin A\), the set \(B\) has five points. Let
\[
\Teich(B)=\Teich(\mathbb{S}^2,B),\qquad \Teich(A)=\Teich(\mathbb{S}^2,A)
\]
be the Teichm\"uller spaces of the sphere with the indicated marked points, and let
\[
f^*:\Teich(B)\to \Teich(A),\qquad i_*:\Teich(B)\to \Teich(A)
\]
be respectively the pullback map induced by \(f\) and the forgetful map induced by the inclusion
\(A\subset B\). The deformation space is
\[
\mathcal D_f:=\{\tau\in \Teich(B): f^*(\tau)=i_*(\tau)\}.
\]
It is a one-dimensional complex submanifold of \(\Teich(B)\); see
\href{https://ar5iv.labs.arxiv.org/html/1606.00723}{[ar5iv.labs.arxiv.org]} and
\href{https://ar5iv.labs.arxiv.org/html/1602.07378}{[ar5iv.labs.arxiv.org]}.

Let \(\Mod(\mathbb{S}^2,B)\) be the pure mapping class group of the five-punctured sphere,
and let \(\Mod(\mathbb{S}^2,A)\) be the pure mapping class group of the four-punctured sphere.

\begin{definition}
Define \(\gamma\) to be the unique, up to isotopy, simple closed curve in
\(\mathbb{S}^2\setminus B\) which bounds an open disk \(D\) such that
\[
D\cap B=\{a_0,v\}=\{c_1,v\}.
\]
Equivalently, \(\gamma\) separates the two marked points \(c_1\) and \(v\) from the other three
marked points \(a_1,a_2,a_3\).
\end{definition}

We prove the following result.

\begin{theorem}
Let \(T_\gamma\in \Mod(\mathbb{S}^2,B)\) be the Dehn twist about the curve \(\gamma\) defined above.
Then:
\begin{enumerate}
\item \(T_\gamma^2\) preserves the deformation space \(\mathcal D_f\).
\item The action of \(T_\gamma^2\) on the set \(\pi_0(\mathcal D_f)\) of connected components of
\(\mathcal D_f\) has infinite orbits.
\end{enumerate}
\end{theorem}

\begin{proof}
We prove the two assertions separately.

\subsection*{1. The Dehn twist \(T_\gamma^2\) preserves \(\mathcal D_f\)}

We use the following standard lifting lemma for Dehn twists. A proof is given in
\href{https://ar5iv.labs.arxiv.org/html/1606.00723}{[ar5iv.labs.arxiv.org]}.

\begin{lemma}[Lifting of Dehn twists]
Let \(f:(\mathbb{S}^2,A)\to(\mathbb{S}^2,B)\) be a branched cover, and let
\(\gamma\) be a simple closed curve in \(\mathbb{S}^2\setminus B\). Let
\(\beta_1,\dots,\beta_k\) be the connected components of \(f^{-1}(\gamma)\), and let
\(d_i\) be the degree of \(f|_{\beta_i}:\beta_i\to\gamma\). If \(m\) is divisible by every
\(d_i\), then \(T_\gamma^m\) lifts under \(f\), and
\[
f^*(T_\gamma^m)=\prod_{i=1}^k T_{\beta_i}^{m/d_i}
\]
in \(\Mod(\mathbb{S}^2,A)\).
\end{lemma}

We also use the following elementary fact about punctured spheres.

\begin{lemma}[Twists around once-punctured disks]
If \(\delta\) is a simple closed curve in \(\mathbb{S}^2\setminus X\), where \(X\) is a finite
marked set, and if \(\delta\) bounds a disk containing at most one point of \(X\), then
\(T_\delta=1\) in \(\Mod(\mathbb{S}^2,X)\).
\end{lemma}

\begin{proof}
A disk with at most one interior marked point and fixed boundary has trivial mapping class
group. Hence the Dehn twist about its boundary is isotopic to the identity relative to the
marked point.
\end{proof}

Now consider our curve \(\gamma\). Let \(D\) be the disk bounded by \(\gamma\) with
\[
D\cap B=\{a_0,v\}.
\]
After forgetting the point \(v\), the disk \(D\) contains exactly the one marked point
\(a_0\in A\). Therefore, by the second lemma,
\[
i_*(T_\gamma)=1
\]
in \(\Mod(\mathbb{S}^2,A)\), and hence
\begin{equation}
i_*(T_\gamma^2)=1. \label{eq:forgetful}
\end{equation}

Next we compute the lift of \(T_\gamma^2\) under \(f\). The disk \(D\) contains exactly one
critical value of \(f\), namely \(v=f(c_2)\). The other critical value is
\(a_1=f(c_1)\), and \(a_1\notin D\). Hence
\[
\Delta:=f^{-1}(D)
\]
is a closed disk, and
\[
f|_{\Delta}:\Delta\to D
\]
is a branched cover of degree \(2\), branched only at \(c_2\) over \(v\). Let
\[
\beta:=\partial\Delta=f^{-1}(\gamma).
\]
Then \(\beta\) is a simple closed curve, and
\[
\deg(f|_{\beta}:\beta\to\gamma)=2.
\]
Thus, by the lifting lemma,
\begin{equation}
f^*(T_\gamma^2)=T_\beta. \label{eq:lift}
\end{equation}

It remains to identify \(T_\beta\) in \(\Mod(\mathbb{S}^2,A)\). The disk \(\Delta\) contains
\(c_2\) and the full preimage \(f^{-1}(a_0)\). Since \(a_0\) is not a critical value,
\(f^{-1}(a_0)\) consists of two distinct points. One of them is \(a_3\), because
\(f(a_3)=a_0\). No other point of \(A\) maps to \(a_0\), because \(A\) is the cycle
\[
a_0\mapsto a_1\mapsto a_2\mapsto a_3\mapsto a_0.
\]
Also \(c_2\notin A\), since by assumption \(c_1\) is the only critical point in the
superattracting \(4\)-cycle. Therefore
\[
\Delta\cap A=\{a_3\}.
\]
Thus \(\beta\) bounds a disk containing exactly one point of \(A\). By the second lemma,
\[
T_\beta=1
\]
in \(\Mod(\mathbb{S}^2,A)\). Combining this with \eqref{eq:lift} gives
\begin{equation}
f^*(T_\gamma^2)=1. \label{eq:flift}
\end{equation}

Now let \(\tau\in\mathcal D_f\). Then \(f^*(\tau)=i_*(\tau)\). Using \eqref{eq:forgetful}
and \eqref{eq:flift}, and the naturality of the mapping class group actions on Teichm\"uller
spaces, we obtain
\[
f^*(T_\gamma^2\cdot \tau)
=
f^*(T_\gamma^2)\cdot f^*(\tau)
=
f^*(\tau)
=
i_*(\tau)
=
i_*(T_\gamma^2)\cdot i_*(\tau)
=
i_*(T_\gamma^2\cdot \tau).
\]
Hence \(T_\gamma^2\cdot\tau\in\mathcal D_f\). Since \(\tau\in\mathcal D_f\) was arbitrary,
\(T_\gamma^2\) preserves \(\mathcal D_f\).

\subsection*{2. The action on components has infinite orbits}

Let
\[
S_f:=\{g\in \Mod(\mathbb{S}^2,B): g(\mathcal D_f)=\mathcal D_f\}
\]
be the stabilizer of \(\mathcal D_f\). Fix a connected component \(C_0\) of
\(\mathcal D_f\), and let
\[
E_f:=\{g\in S_f: g(C_0)=C_0\}
\]
be the stabilizer of that component. We have just proved that \(T_\gamma^2\in S_f\).

The following sharper statement is proved in the course of the proof of the main theorem of
\href{https://ar5iv.labs.arxiv.org/html/1602.07378}{[ar5iv.labs.arxiv.org]}.

\begin{proposition}[{\cite[proof of Theorem 1.3]{disconnected}}]
For the quadratic rational maps considered here, the Dehn twist \(T_\gamma^2\) has infinite
order in the quotient \(S_f/E_f\). Equivalently,
\[
(T_\gamma^2)^n\notin E_f
\]
for every \(n\neq 0\).
\end{proposition}

\begin{proof}[Proof of the proposition as given in the cited paper]
The proof in
\href{https://ar5iv.labs.arxiv.org/html/1602.07378}{[ar5iv.labs.arxiv.org]} compares the
stabilizer \(S_f\) of the whole deformation space with the stabilizer \(E_f\) of a connected
component. The authors produce an element whose no nontrivial power lies in \(E_f\); in the
notation of the present paper, that element is exactly \(T_\gamma^2\). This shows that the
cosets
\[
(T_\gamma^2)^n E_f\qquad (n\in\mathbb Z)
\]
are all distinct, and hence the orbit of \(C_0\) under the cyclic group generated by
\(T_\gamma^2\) is infinite.
\end{proof}

Therefore, for every integer \(n\neq 0\),
\[
(T_\gamma^2)^n(C_0)\neq C_0.
\]
Thus the components
\[
(T_\gamma^2)^n(C_0)\qquad (n\in\mathbb Z)
\]
are pairwise distinct. In particular, the action of \(T_\gamma^2\) on
\(\pi_0(\mathcal D_f)\) has an infinite orbit.
\end{proof}

\begin{remark}
The assumption that \(v=f(c_2)\) is not in the cycle ensures that \(c_2\notin A\) and that
\(B=A\cup\{v\}\) consists of five distinct marked points. This is precisely the class
\(\operatorname{Per}_4(0)^*\) of quadratic rational maps with a superattracting \(4\)-cycle
containing exactly one critical point; see
\href{https://ar5iv.labs.arxiv.org/html/1602.07378}{[ar5iv.labs.arxiv.org]}.
\end{remark}

\begin{thebibliography}{9}

\bibitem{defs}
\href{https://ar5iv.labs.arxiv.org/html/1606.00723}{[ar5iv.labs.arxiv.org]}
\emph{Deformation spaces and quadratic rational maps}, arXiv:1606.00723.

\bibitem{disconnected}
\href{https://ar5iv.labs.arxiv.org/html/1602.07378}{[ar5iv.labs.arxiv.org]}
\emph{A disconnected deformation space of rational maps}, arXiv:1602.07378;
published version: \emph{J. Mod. Dyn.} \textbf{11} (2017), 409--425.

\bibitem{FKS}
\href{https://www.math.ksu.edu/~tanyaf/FKS.pdf}{[math.ksu.edu]}
\emph{On deformation spaces of quadratic rational functions}.

\end{thebibliography}

\end{document}
```